\documentclass[11pt,a4paper]{article}
\usepackage[utf8]{inputenc}
\usepackage[T1]{fontenc}
\usepackage{geometry}
\usepackage{xcolor}
\usepackage{hyperref}
\usepackage{titlesec}

\geometry{margin=1in}
\definecolor{tumblue}{RGB}{0,101,189}

\hypersetup{
    colorlinks=true,
    linkcolor=tumblue,
    urlcolor=tumblue,
    pdftitle={Cover Letter - Ismail AISSAOUI},
}

\begin{document}

\noindent \textbf{Ismail AISSAOUI} \\
State Engineer in Artificial Intelligence (ESI-SBA) \\
Email: ismail.aissaoui@example.com \\
Phone: [Your Phone Number] \\
LinkedIn: [Your Link]

\bigskip

\begin{flushright}
    Professor Andreas Herkersdorf \\
    \textbf{Technical University of Munich (TUM)} \\
    Chair of Integrated Systems (LIS) \\
    Munich, Germany \\
    \medskip
    May 5, 2026
\end{flushright}

\bigskip

\noindent \textbf{Subject: Application for Doctoral Research - Hardware-Aware Edge AI}

\bigskip

Dear Professor Herkersdorf,

I am writing to formally submit my application for a doctoral position under your supervision at the Chair of Integrated Systems at the Technical University of Munich. As an AI Engineer with a focus on the hardware-aware deployment of complex neural architectures, I am deeply interested in your work on HW/SW co-design and application-specific AI accelerators.

Currently, I am concluding my state engineering thesis at **Sonatrach**, where I have been developing a **Generative Digital Twin** for industrial health monitoring. A primary challenge of this project is the real-time processing of high-frequency sensor data on industrial edge devices. To address this, I have been working with **State Space Models (Mamba-SSM)** and **Graph Neural Networks (GNNs)**, exploring their computational efficiency compared to standard Transformer models.

I am particularly interested in the research activities at LIS regarding the optimization of AI models for embedded and edge systems. My technical background in Python and PyTorch, combined with my practical experience in deploying models for industrial monitoring, allows me to bridge the gap between high-level AI research and low-level hardware constraints. I am eager to research how the linear recurrence and specialized hardware requirements of **State Space Models** can be co-optimized with custom hardware accelerators (FPGA/ASIC) to enable the next generation of autonomous industrial monitoring.

TUM is my top choice for a PhD due to its leadership in both hardware design and artificial intelligence. I am a highly motivated researcher, and I believe my background from ESI-SBA (recognized as H+ in Germany) and my practical industrial work at Sonatrach have prepared me for the interdisciplinary challenges of your group.

I have attached my academic transcripts, my CV, and a research proposal for your review. I would be honored to discuss how my technical expertise can contribute to your laboratory’s mission.

Thank you for your time and consideration.

Sincerely,

\bigskip

\noindent Ismail AISSAOUI

\end{document}
